\section{Closeness of low-energy subspaces under a perturbation}

Let $H$ and $H'$ be self-adjoint operators with orthonormal eigenbases
$(b_j)$ and $(b'_j)$ and eigenvalues $a_j$ and $a'_j$. For an energy $c$
write $T(c)$ for the span of the $b_j$ with $a_j\le c$, and $T'(c)$ for the
span of the $b'_j$ with $a'_j\le c$.

\begin{definition}[Low-energy subspace]
    \label{def:spectral_span_le}
    \lean{Submodule.spectralSpanLE}
    \leanok
    The subspace $T(c)$ is the span of the eigenvectors $b_j$ with
    eigenvalue $a_j\le c$.
\end{definition}

\begin{theorem}[Markov's inequality]
    \label{thm:spectral_span_markov}
    \lean{Submodule.mul_norm_starProjection_orthogonal_spectralSpanLE_le}
    \leanok
    \uses{def:spectral_span_le}
    If every eigenvalue $a_j>c$ satisfies $|a_j-\mu|\ge D\ge0$, then the
    projection $P^\perp$ onto the orthogonal complement of $T(c)$ satisfies
    $D\|P^\perp x\|\le\|(H-\mu)x\|$ for every $x$.
\end{theorem}

\begin{proof}
    \leanok
    Expand $x=\sum_jx_jb_j$. Then
    $\|(H-\mu)x\|^2=\sum_j|a_j-\mu|^2|x_j|^2\ge D^2\sum_{a_j>c}|x_j|^2
        =D^2\|P^\perp x\|^2$.
\end{proof}

\begin{theorem}[Eigenvalue comparison from a low-energy residual]
    \label{thm:low_spectral_eigenvalue_comparison}
    \lean{Matrix.IsHermitian.eigenvalues₀_le_add_of_norm_domRestrict_low_le}
    \notready
    \uses{def:spectral_span_le}
    Let $H,F\in\MN{d}$ be Hermitian, with eigenvalues listed in decreasing
    order, and let $G$ be the closed spectral subspace of $F$ at energy $c$.
    Suppose $\|(H-F)|_G\|\le\xi$. For every $0\le i<d$ with
    $\lambda_i(F)\le c$, one has $\lambda_i(H)\le\lambda_i(F)+\xi$.
    This is a variational estimate for use with the restricted perturbation
    bounds in~\cite[lines 955--968]{Arad2016RigorousRG}.
\end{theorem}
\begin{proof}
    \notready
    The first $i+1$ eigenvectors of $H$ and the last $d-i$ eigenvectors of
    $F$ span subspaces with a common nonzero vector $x$. The latter subspace
    is contained in $G$. Spectral expansion and the residual estimate give
    \begin{align}
        \lambda_i(H)\|x\|^2
         & \le \operatorname{Re}\langle Hx,x\rangle
        \le \operatorname{Re}\langle Fx,x\rangle+\xi\|x\|^2
        \le (\lambda_i(F)+\xi)\|x\|^2.
        \notag
    \end{align}
    Divide by $\|x\|^2>0$.
\end{proof}

\begin{theorem}[One window]
    \label{thm:spectral_span_window}
    \lean{Submodule.mul_norm_starProjection_orthogonal_spectralSpanLE_le_of_window}
    \leanok
    \uses{thm:spectral_span_markov}
    Suppose $\|(H-H')\varphi\|\le\xi\|\varphi\|$ for every $\varphi\in T(c)$,
    and let $h,D\ge0$ with $\mu+D\le c'$. If $\psi\in T(c)$ is supported on
    eigenvectors $b_j$ with $|a_j-\mu|\le h$, then the projection $P'^\perp$
    onto the orthogonal complement of $T'(c')$ satisfies
    $D\|P'^\perp\psi\|\le(\xi+h)\|\psi\|$.
\end{theorem}

\begin{proof}
    \leanok
    Every eigenvalue $a'_k>c'$ satisfies $|a'_k-\mu|\ge D$. By the previous
    theorem,
    $D\|P'^\perp\psi\|\le\|(H'-\mu)\psi\|\le\|(H'-H)\psi\|
        +\|(H-\mu)\psi\|\le(\xi+h)\|\psi\|$.
\end{proof}

\begin{theorem}[Closeness of the low-energy subspaces]
    \label{thm:spectral_span_close_of_perturbation}
    \lean{Submodule.isCloseTo_spectralSpanLE_of_perturbation}
    \leanok
    \uses{thm:spectral_span_window, def:subspace_close}
    Suppose $\|(H-H')\varphi\|\le\xi\|\varphi\|$ for every $\varphi\in T(c)$,
    every eigenvalue $a_j\le c$ is at least $m$, $m<c\le c'$ and $h>0$. Let
    $N=\lceil(c-m)/(2h)\rceil$. If $N(\xi+h)^2\le\delta(c'-c+h)^2$, then
    $T'(c')$ is $\delta$-close to $T(c)$. This is the closeness step in the
    proof of the second item of Lemma lem:hard-trunc
    of~\cite[lines 953--969]{Arad2016RigorousRG}, with the number of windows
    corrected~\cite{gap:alvv16_hard_truncation_windows}.
\end{theorem}

\begin{proof}
    \leanok
    Cover $[m,c]$ by the $N$ closed windows of width $2h$ centred at
    $c-h-2hw$ for $w<N$. Split $\varphi\in T(c)$ as
    $\varphi=\sum_{w<N}\psi_w$, where $\psi_w$ keeps the coefficients of the
    $b_j$ with $\lceil(c-a_j)/(2h)\rceil-1=w$, counting $a_j=c$ in the top
    window. Each $\psi_w$ lies in $T(c)$ and is supported within $h$ of its
    centre, and every centre plus $c'-c+h$ is at most $c'$. By the previous
    theorem and the Cauchy--Schwarz inequality,
    \[
        (c'-c+h)\|P'^\perp\varphi\|\le(\xi+h)\sum_w\|\psi_w\|
        \le(\xi+h)\sqrt{N}\,\|\varphi\|,
    \]
    so $\|P'^\perp\varphi\|^2\le\delta\|\varphi\|^2$ and
    $\|P'\varphi\|^2\ge(1-\delta)\|\varphi\|^2$.
\end{proof}

\begin{theorem}[A low-energy cluster from restricted residual estimates]
    \label{thm:low_spectral_cluster_restricted_residuals}
    \lean{Submodule.low_spectral_cluster_of_restricted_residuals}
    \leanok
    \uses{def:spectral_span_le, def:subspace_close}
    Let $H,H'$ act on a finite-dimensional complex inner-product space,
    with complete orthonormal eigenbases and real eigenvalues $a_j,a'_j$.
    Suppose $\gamma>0$, $0\le\xi\le\gamma/10^6$, and every $a_j$ equals
    $e$ or is at least $e+\gamma$. Assume the two restricted operator
    estimates
    \begin{align}
        \bigl\|(H-H')|_{T(e+\gamma/40)}\bigr\|    & \le\xi,
        \label{eq:low_cluster_original_residual}            \\
        \bigl\|(H-H')|_{T'(e+19\gamma/20)}\bigr\| & \le\xi.
        \label{eq:low_cluster_perturbed_residual}
    \end{align}
    Put $G=T(e)$, $U=T'(e+\gamma/20)$ and $V=T'(e+19\gamma/20)$.
    Then $U=V$ and $\dim U=\dim G$. Each of $U$ and $G$ is
    $1/200$-close to the other. For every
    $e+\gamma/20\le\tau\le e+19\gamma/20$, the space $T'(\tau)$ equals
    $U$ and retains both approximation bounds. Every perturbed eigenvalue
    satisfies
    \begin{align}
        a'_j\le e+\gamma/10
        \quad\text{or}\quad e+9\gamma/10\le a'_j.
        \label{eq:low_cluster_eigenvalue_gap}
    \end{align}
    The cutoffs are absolute energies relative to $e$; no comparison
    with the perturbed minimum energy is asserted. This is a
    finite-dimensional consequence of the estimates
    in~\cite[lines 955--968 and 1010--1024]{Arad2016RigorousRG}, using the
    corrected window count~\cite{gap:alvv16_hard_truncation_windows}.
\end{theorem}

\begin{proof}
    \leanok
    \uses{thm:spectral_span_close_of_perturbation,
        thm:subspace_close_nested_equal, thm:subspace_close_dimension,
        thm:spectral_span_quadratic_form}
    The original spectral dichotomy gives $H\ge eI$ and
    $T(e+\gamma/40)=T(e+39\gamma/40)=G$. By
    (\ref{eq:low_cluster_perturbed_residual}), a unit eigenvector $x$ of
    $H'$ with eigenvalue $a'_j\le e+19\gamma/20$ satisfies
    $e\le\operatorname{Re}\langle x,Hx\rangle\le a'_j+\xi$.
    Thus these perturbed eigenvalues are at least $e-\xi$.

    Apply the low-energy perturbation estimate with
    $c=e+\gamma/40$, $c'=e+\gamma/20$, $m=e$ and then with the two
    operators interchanged and
    $c=e+19\gamma/20$, $c'=e+39\gamma/40$, $m=e-\xi$.
    In each application take $h=\gamma/10^6$. The source width is at
    most $\gamma$, the cutoff separation is $\gamma/40$, and the number
    of windows is at most $500001$. Consequently
    \begin{align}
        N(\xi+h)^2
         & \le 500001(2\gamma/10^6)^2
        \le \frac1{200}(\gamma/40)^2
        \le \frac1{200}(c'-c+h)^2.
        \label{eq:low_cluster_window_bound}
    \end{align}
    The two applications show that $U$ is $1/200$-close to $G$ and $G$
    is $1/200$-close to $V$. Since $U\subseteq V$, the nested-subspace
    theorem gives $U=V$, equality of dimensions with $G$, and both
    approximation bounds for every intermediate cutoff. If
    $e+\gamma/20<a'_j\le e+19\gamma/20$, its eigenvector belongs to
    $V=U$, whereas its own coordinate must vanish in $U$. This
    contradiction proves~(\ref{eq:low_cluster_eigenvalue_gap}).
\end{proof}


\begin{theorem}[Energy on a low spectral subspace]
    \label{thm:spectral_span_quadratic_form}
    \lean{OrthonormalBasis.re_inner_apply_eq_sum_of_eigen,
        Submodule.re_inner_apply_le_of_mem_spectralSpanLE}
    \leanok
    \uses{def:spectral_span_le}
    Let $(b_j)$ be a finite orthonormal basis of a complex inner-product
    space, and let $H$ be a linear operator satisfying $Hb_j=a_jb_j$ with
    $a_j\in\R$. For every $x$,
    \begin{align}
        \operatorname{Re}\langle x,Hx\rangle
         & =\sum_j a_j|\langle b_j,x\rangle|^2.
        \label{eq:spectral_span_energy_coordinates}
    \end{align}
    In particular, if $x\in T(c)$, then
    $\operatorname{Re}\langle x,Hx\rangle\le c\|x\|^2$.
\end{theorem}
\begin{proof}
    \leanok
    Express $x$ and $Hx$ in the orthonormal basis to obtain
    (\ref{eq:spectral_span_energy_coordinates}). For $x\in T(c)$, all
    coordinates with $a_j>c$ vanish. Bound the remaining coefficients by
    $c$ and apply Parseval's identity.
\end{proof}

\begin{theorem}[The low spectral projection commutes with the operator]
    \label{thm:spectral_span_projection_commutes}
    \lean{Submodule.commute_starProjection_spectralSpanLE}
    \leanok
    \uses{def:spectral_span_le}
    Under the same eigenbasis assumptions, the orthogonal projection
    onto $T(c)$ commutes with $H$, for every real $c$.
\end{theorem}
\begin{proof}
    \leanok
    In that basis the complementary projection retains precisely the
    coordinates with $a_j>c$, while $H$ multiplies each coordinate by $a_j$.
    These two operations commute, and so do $H$ and the low projection.
    These are spectral consequences used
    in~\cite[lines 955--968 and 1010--1024]{Arad2016RigorousRG}.
\end{proof}

\begin{theorem}[Relative low-energy spaces above the actual ground energies]
    \label{thm:relative_low_energy_spaces_close}
    \lean{Matrix.IsHermitian.groundEnergy_compare_and_isCloseTo_spectralSpanLE_of_norm_domRestrict_le}
    \notready
    \uses{def:spectral_span_le,def:subspace_close}
    Let $H,F$ be Hermitian matrices on a nonzero finite-dimensional space,
    with $F\le H$. Let $E$ and $E'$ be their respective least eigenvalues,
    and write $T_H(c)$ and $T_F(c)$ for their closed spectral subspaces
    at the absolute energy $c$.
    Suppose $\eta,\xi,\delta>0$ and
    \begin{align}
        \|(H-F)|_{T_H(E+\eta)}\| & \le\xi, &
        \|(H-F)|_{T_F(E+\eta)}\| & \le\xi.
        \label{eq:relative_low_energy_residuals}
    \end{align}
    Put $N=\lceil\eta/(2\xi)\rceil$ and
    $\eta'=\eta+\xi+2\xi\sqrt{N/\delta}$.
    Then $E'\le E\le E'+\xi$,
    $T_F(E'+\eta')$ is $\delta$-close to $T_H(E+\eta)$, and
    $T_H(E+\eta')$ is $\delta$-close to $T_F(E'+\eta)$.
    Both relative cutoffs use the actual respective ground energy.
    This is the finite-dimensional relative-energy consequence of
    \cite[lines 955--968]{Arad2016RigorousRG}, with the window count and
    ground-energy shift corrected in
    \cite{gap:alvv16_hard_truncation_windows}.
\end{theorem}
\begin{proof}
    \notready
    \uses{thm:low_spectral_eigenvalue_comparison,
        thm:spectral_span_close_of_perturbation}
    Weyl monotonicity gives $E'\le E$. Since the lowest eigenvalue of
    $F$ lies below $E+\eta$, the restricted eigenvalue comparison gives
    $E\le E'+\xi$.
    Apply the window estimate with width $h=\xi$ to the spaces of $H$
    between $E$ and $E+\eta$. For the enlarged target cutoff, the
    denominator satisfies
    \begin{align}
        (E'+\eta')-(E+\eta)+\xi
         & \ge\xi+2\xi\sqrt{N/\delta}
        \ge2\xi\sqrt{N/\delta}.
        \notag
    \end{align}
    Its squared value times $\delta$ is at least $N(2\xi)^2$, so the
    finite-window budget holds. The target cutoff is at least the source
    cutoff by the same ground-energy comparison.
    For the reverse direction, $E'\le E$ gives
    $T_F(E'+\eta)\subseteq T_F(E+\eta)$, so the second residual estimate
    remains valid on the smaller space. Its spectral width is again
    $\eta$, and the reverse target denominator is larger still.
    Apply the same window argument with $H,F$ interchanged.
\end{proof}
